\documentclass[a4paper,twoside]{article}

\usepackage{morewrites}

\usepackage[svgnames,dvipsnames,x11names]{xcolor}
\usepackage{graphicx}
\usepackage{texsmith-lists}
\usepackage[english]{babel}
\usepackage[colorlinks=true, linkcolor=NavyBlue, urlcolor=NavyBlue, citecolor=NavyBlue]{hyperref}
\usepackage[a4paper]{geometry}
\usepackage{ts-fonts}
\usepackage{float}
\usepackage{graphicx}
\usepackage{amsmath}
\usepackage{seqsplit}
\usepackage{ts-code}

\usepackage{lastpage}
\usepackage{refcount}
\makeatletter
\def\ps@plain{
\let\@oddhead\@empty
\let\@evenhead\@empty
\def\@oddfoot{
\hfil
\ifnum\getpagerefnumber{LastPage}>1\relax
\thepage
\fi
\hfil}
\let\@evenfoot\@oddfoot
}
\makeatother

\title{Book}
\author{}\date{}
\begin{document}
\maketitle
\phantomsection\label{einstein}
For this example we want to render homage to Albert Einstein by creating a small book using the \href{https://en.wikipedia.org/wiki/Albert\_Einstein}{Wikipedia}
as our source.

The text was converted to \LaTeX{} then built using the \tscodeinline{book} template.

This example demonstrates the use of front matter for citations, abbreviations, glossary, and index.

For this example, we decided to use a small paper size (A5) and a sans serif font (Adventor). The book is structured in parts and specific slots are used for the colophon, dedication, and preface. Here is the front matter used for this example:
\begin{tscode}[lang=yaml, engine=pygments]
\begin{Verbatim}[commandchars=\\\{\},breaklines, breakanywhere, commandchars=\\\{\}]
\PY{n+nt}{title}\PY{p}{:}\PY{+w}{ }\PY{l+lScalar+lScalarPlain}{Albert}\PY{l+lScalar+lScalarPlain}{ }\PY{l+lScalar+lScalarPlain}{Einstein}
\PY{n+nt}{subtitle}\PY{p}{:}\PY{+w}{ }\PY{l+lScalar+lScalarPlain}{His}\PY{l+lScalar+lScalarPlain}{ }\PY{l+lScalar+lScalarPlain}{Life}\PY{l+lScalar+lScalarPlain}{ }\PY{l+lScalar+lScalarPlain}{and}\PY{l+lScalar+lScalarPlain}{ }\PY{l+lScalar+lScalarPlain}{Achievements}
\PY{n+nt}{date}\PY{p}{:}\PY{+w}{ }\PY{l+lScalar+lScalarPlain}{2025\PYZhy{}03\PYZhy{}15}
\PY{n+nt}{edition}\PY{p}{:}\PY{+w}{ }\PY{l+lScalar+lScalarPlain}{Wikipedia}\PY{l+lScalar+lScalarPlain}{ }\PY{l+lScalar+lScalarPlain}{Adaptation}
\PY{n+nt}{author}\PY{p}{:}\PY{+w}{ }\PY{l+lScalar+lScalarPlain}{Wikipedia}\PY{l+lScalar+lScalarPlain}{ }\PY{l+lScalar+lScalarPlain}{Contributors}
\PY{n+nt}{publisher}\PY{p}{:}\PY{+w}{ }\PY{l+lScalar+lScalarPlain}{Wikipedia}\PY{l+lScalar+lScalarPlain}{ }\PY{l+lScalar+lScalarPlain}{Sourcebooks}
\PY{n+nt}{imprint}\PY{p}{:}
\PY{+w}{  }\PY{n+nt}{thanks}\PY{p}{:}\PY{+w}{ }\PY{p+pIndicator}{|}
\PY{+w}{    }\PY{n+no}{Adapted from the Wikipedia article “Albert Einstein” (retrieved ).}
\PY{+w}{    }\PY{n+no}{Thank you to the Wikipedia contributors whose work makes this example}
\PY{+w}{    }\PY{n+no}{possible.}
\PY{+w}{  }\PY{n+nt}{copyright}\PY{p}{:}\PY{+w}{ }\PY{p+pIndicator}{|}
\PY{+w}{    }\PY{n+no}{Wikipedia contributors. Text licensed under Creative Commons Attribution\PYZhy{}ShareAlike}
\PY{+w}{    }\PY{n+no}{International (CC BY\PYZhy{}SA).}
\PY{+w}{  }\PY{n+nt}{license}\PY{p}{:}\PY{+w}{ }\PY{p+pIndicator}{|}
\PY{+w}{    }\PY{n+no}{This TeXSmith example is derived from Wikipedia and is not affiliated}
\PY{+w}{    }\PY{n+no}{with TeXSmith. Reuse must credit Wikipedia and share under CC BY\PYZhy{}SA .}
\PY{n+nt}{press}\PY{p}{:}
\PY{+w}{  }\PY{n+nt}{paper}\PY{p}{:}\PY{+w}{ }\PY{l+lScalar+lScalarPlain}{a5}
\PY{+w}{  }\PY{n+nt}{template}\PY{p}{:}\PY{+w}{ }\PY{l+lScalar+lScalarPlain}{book}
\PY{+w}{  }\PY{n+nt}{base\PYZus{}level}\PY{p}{:}\PY{+w}{ }\PY{l+lScalar+lScalarPlain}{part}
\PY{+w}{  }\PY{n+nt}{fonts}\PY{p}{:}\PY{+w}{ }\PY{l+lScalar+lScalarPlain}{adventor}
\PY{+w}{  }\PY{n+nt}{callouts}\PY{p}{:}
\PY{+w}{    }\PY{n+nt}{style}\PY{p}{:}\PY{+w}{ }\PY{l+lScalar+lScalarPlain}{classic}
\PY{+w}{  }\PY{n+nt}{slots}\PY{p}{:}
\PY{+w}{    }\PY{n+nt}{colophon}\PY{p}{:}\PY{+w}{ }\PY{l+lScalar+lScalarPlain}{Colophon}
\PY{+w}{    }\PY{n+nt}{dedication}\PY{p}{:}\PY{+w}{ }\PY{l+lScalar+lScalarPlain}{Dedication}
\PY{+w}{    }\PY{n+nt}{preface}\PY{p}{:}\PY{+w}{ }\PY{l+lScalar+lScalarPlain}{Preface}
\end{Verbatim}
\end{tscode}
\begin{figure}[H]
\centering
\includegraphics[width=\linewidth]{snippet-<HASH>.pdf}
\end{figure}
The example can be built independently using the CLI in the \tscodeinline{examples/book/} folder. The engine can be chosen between \tscodeinline{tectonic}, \tscodeinline{xelatex}, and \tscodeinline{lualatex}
as follows:
\begin{tscode}[lang=bash, engine=pygments]
\begin{Verbatim}[commandchars=\\\{\},breaklines, breakanywhere, commandchars=\\\{\}]
texsmith\PY{+w}{ }book.md\PY{+w}{ }<STEM>.bib\PY{+w}{ }\PYZhy{}\PYZhy{}template\PY{+w}{ }book\PY{+w}{ }\PYZhy{}\PYZhy{}build\PY{+w}{ }\PYZhy{}\PYZhy{}engine\PY{+w}{ }tectonic
\end{Verbatim}
\end{tscode}
The power of TeXSmith is that it can generate such complex documents from simple Markdown with no \LaTeX{} knowledge required and no \LaTeX{} toolchain installed. Fonts, toolchain, and image conversion are all handled automatically by TeXSmith in a single command.
\end{document}
